\chapter{Arquitetura de Software}
\label{cap:arquitetura}

\begin{edital}
Arquitetura de software; interoperabilidade; arquitetura e linguagem
orientada a serviços; web services; mensageria; API, Swagger; arquitetura
orientada a objetos; aplicações para ambiente web; servidor de aplicações
$\times$ servidor web; internet/extranet/intranet/portal; arquitetura
hexagonal; microsserviços (orquestração e API gateway); containers;
transações distribuídas.
\end{edital}
\peso{ALTO --- nuvem, REST$\times$SOAP, microsserviços e containers são
recorrentes.}

\section{Servidor web $\times$ servidor de aplicação}

\begin{center}
\begin{tabular}{@{}p{2.4cm}p{4.6cm}p{4.6cm}@{}}
\toprule
& \textbf{Servidor web} & \textbf{Servidor de aplicação} \\
\midrule
Função & atende HTTP, serve conteúdo estático/dinâmico & executa \textbf{lógica de negócio}, integra com BD \\
Exemplos & Apache, NGINX & JBoss/WildFly, WebLogic, Tomcat (parcial) \\
\bottomrule
\end{tabular}
\end{center}

\begin{pegadinha}
A FGV troca os papéis (``servidor web processa lógica de negócio''). O
\textbf{web} recebe a requisição HTTP; o \textbf{de aplicação} roda a regra
de negócio.
\end{pegadinha}

\section{Estilos de arquitetura}

\begin{conceito}[Do monólito aos microsserviços]
\begin{itemize}
\item \textbf{Monolítica:} tudo num artefato único, implantado junto.
Simples de começar; acopla e dificulta escalar partes isoladas.
\item \textbf{Cliente-servidor, N camadas, P2P, barramento de mensagens.}
\item \textbf{SOA (orientada a serviços):} serviços reutilizáveis, contrato
explícito, baixo acoplamento, interoperabilidade. Web services são a
tecnologia comum.
\item \textbf{Microsserviços:} serviços pequenos, autônomos,
\textbf{implantáveis independentemente}, cada um com \textbf{seu próprio
banco} (não compartilham base --- isso reduz acoplamento). Comunicação leve
(REST/mensageria).
\item \textbf{Arquitetura hexagonal (Ports \& Adapters):} separa a
\textbf{lógica de negócio} do mundo externo por \textbf{portas} (interfaces)
e \textbf{adaptadores}; permite trocar UI/BD/serviços sem tocar no núcleo.
\end{itemize}
\end{conceito}

\begin{pegadinha}
``Microsserviços compartilham o mesmo banco'' = \textbf{falso} (aumenta
acoplamento, contra o princípio). ``Monólito não pode ser distribuído'' =
falso. Baixo acoplamento significa que mudança em um serviço \textbf{não}
obriga mudança no outro --- se a alternativa disser que ``se reflete
diretamente'', descarte.
\end{pegadinha}

\section{Integração: REST $\times$ SOAP, Web Services}

\begin{center}
\begin{tabular}{@{}p{2.6cm}p{4.6cm}p{4.6cm}@{}}
\toprule
& \textbf{REST} & \textbf{SOAP} \\
\midrule
Estilo & arquitetural, sobre HTTP & protocolo baseado em XML \\
Formato & JSON (comum), leve & XML (envelope), verboso \\
Contrato & OpenAPI/Swagger & WSDL \\
Estado & \textbf{stateless} & pode ter padrões WS-* \\
Vantagem & leve, escala, independe de plataforma & contratos formais, WS-Security \\
\bottomrule
\end{tabular}
\end{center}

\begin{itemize}
\item \textbf{REST é stateless:} cada requisição carrega tudo; o servidor
não guarda sessão $\to$ escala horizontalmente sem afinidade de sessão. Cache
fica no \textbf{cliente}.
\item \textbf{API Gateway:} ponto único de entrada dos microsserviços
(roteamento, autenticação, rate limit, agregação).
\item \textbf{Métodos HTTP:} GET (ler), POST (criar), \textbf{PUT
(substituir inteiro)}, \textbf{PATCH (atualizar parcial)}, DELETE. Códigos:
200 OK, 201 Created, 204 No Content, 400, 401, 403, 404, 500.
\item \textbf{UDDI:} diretório (legado) para \textbf{descoberta/registro} de
web services SOAP, junto do WSDL (contrato) e SOAP (mensagem). Trio clássico
WS-*: \textbf{SOAP + WSDL + UDDI}.
\item \textbf{Swagger / OpenAPI:} o \textbf{OpenAPI} é a especificação
padrão para \textbf{documentar e contratar APIs REST} (endpoints,
parâmetros, respostas); \textbf{Swagger} é o conjunto de ferramentas (Swagger
UI, editor, codegen) em cima dela. É ``o WSDL do REST''.
\end{itemize}

\begin{pegadinha}
\textbf{PUT $\times$ PATCH} (substituição total $\times$ parcial); dizer que
SOAP não tem contrato formal contradiz o próprio SOAP (que usa WSDL); UDDI é
\textbf{descoberta}, WSDL é \textbf{contrato}, SOAP é a \textbf{mensagem} ---
não misture os três papéis.
\end{pegadinha}

\subsection{Mensageria (comunicação assíncrona)}

Desacopla produtor e consumidor: quem envia não espera quem processa.

\begin{center}
\begin{tabular}{@{}p{3.2cm}p{4.6cm}p{3cm}@{}}
\toprule
\textbf{Modelo} & \textbf{Como funciona} & \textbf{Exemplo} \\
\midrule
Fila (queue) & mensagem consumida por \textbf{um} consumidor (point-to-point) & RabbitMQ, SQS \\
Publish/Subscribe (tópico) & mensagem entregue a \textbf{vários} assinantes & Kafka, SNS \\
\bottomrule
\end{tabular}
\end{center}

\textbf{Apache Kafka:} plataforma de \textbf{streaming} de eventos ---
produtores publicam em \textbf{tópicos} (particionados), consumidores leem
em \textbf{grupos}; guarda o log de eventos (permite reprocessar). Alta
vazão, escalável.

Benefícios da mensageria: desacoplamento, absorção de picos (buffer),
resiliência. O \textbf{broker} é o intermediário (Kafka, RabbitMQ). Casa com
\textbf{microsserviços} e com os padrões \textbf{Saga} (eventos de
compensação) e \textbf{Event Sourcing}.

\begin{pegadinha}
Fila (\textbf{um} consumidor) $\times$ tópico/pub-sub (\textbf{vários}) ---
não inverta. Mensageria é \textbf{assíncrona} (diferente da chamada REST
síncrona).
\end{pegadinha}

\section{Ambientes de rede corporativa}

\begin{center}
\begin{tabular}{@{}p{2.4cm}p{9cm}@{}}
\toprule
\textbf{Termo} & \textbf{Alcance} \\
\midrule
Internet & rede pública global \\
Intranet & rede interna da organização (só funcionários) \\
Extranet & acesso controlado a parceiros/clientes externos autorizados \\
Portal & ponto único que centraliza informação/serviços \\
\bottomrule
\end{tabular}
\end{center}

\begin{pegadinha}
Trocar intranet$\leftrightarrow$extranet$\leftrightarrow$internet. Extranet =
interno \textbf{+} parceiros externos autorizados (o meio-termo).
\end{pegadinha}

\section{Nuvem e escalabilidade}

\begin{itemize}
\item \textbf{Modelos de serviço:} \textbf{IaaS} (infra), \textbf{PaaS}
(plataforma), \textbf{SaaS} (software pronto). Regra: quanto mais ``as a
Service'', menos você gerencia.
\item \textbf{Modelos de implantação:} privada, pública, híbrida.
\item \textbf{Escalabilidade horizontal (scale out):} \textbf{adicionar
instâncias/nós}.
\item \textbf{Escalabilidade vertical (scale up):} \textbf{adicionar
recursos à instância} existente (mais CPU/RAM).
\item \textbf{Elasticidade:} ajustar capacidade automaticamente conforme a
demanda. \textbf{Cloud bursting:} estende para nuvem pública no pico de
demanda.
\item \textbf{Containers (Docker) $\times$ VM:} container compartilha o
\textbf{kernel} do SO (leve, rápido); VM tem SO convidado completo sobre um
hipervisor (isola mais, pesa mais). \textbf{Kubernetes} orquestra containers.
\item \textbf{Hipervisor Tipo 1} (bare-metal, roda direto no hardware: ESXi,
Hyper-V, KVM) $\times$ \textbf{Tipo 2} (hosted, roda sobre um SO: VirtualBox,
VMware Workstation).
\item \textbf{Virtualização total} (SO convidado não sabe que é virtual)
$\times$ \textbf{paravirtualização} (SO convidado é modificado e coopera com
o hipervisor --- mais rápido). \textbf{VDI} = virtualização de desktops
entregues remotamente.
\end{itemize}

\begin{pegadinha}
``Adicionar recursos à instância'' é sempre \textbf{vertical}, jamais
horizontal. Container $\neq$ VM (kernel compartilhado $\times$ SO próprio). Em
questões de custo-benefício, some a capacidade das novas instâncias e compare
o preço antes de decidir entre escalar vertical ou horizontal.
\end{pegadinha}

\section{Transações distribuídas}

\begin{itemize}
\item \textbf{2PC (Two-Phase Commit):} coordenador + participantes; exige
\textbf{unanimidade} (todos confirmam) para commit. Bloqueante.
\item \textbf{Saga:} sequência de transações locais com \textbf{compensação}
em caso de falha (padrão para microsserviços).
\end{itemize}

\begin{jacaiu}
Servidor web $\times$ de aplicação; SOA e web services (baixo acoplamento,
REST); arquitetura hexagonal $\times$ microsserviços (não compartilham
banco); internet/extranet/intranet/portal; escalabilidade horizontal $\times$
vertical (com cálculo de custo-benefício); REST stateless; container $\times$
VM; API gateway; taints/tolerations no Kubernetes; cloud bursting; 2PC
$\times$ Raft. Rode \texttt{./quiz.py arquitetura}.
\end{jacaiu}

\section{Alta probabilidade / pesquisa extra}

\begin{itemize}
\item \textbf{12-Factor App} (boas práticas de app nativa de nuvem).
\item \textbf{DDD (Domain-Driven Design):} \textbf{Aggregate} (fronteira de
consistência, protege invariantes --- não expõe referência a cada entidade
interna), \textbf{Repository} (coleção/persistência do agregado, não instancia
por métodos externos), \textbf{Factory} (criação), \textbf{Ubiquitous
Language} (vocabulário comum domínio$\leftrightarrow$código), eventos de
domínio \textbf{imutáveis} (fato passado). Casou com microsserviços no TJ-RJ e
no MPU.
\item \textbf{Service mesh} (Istio) $\times$ API gateway: mesh cuida da
comunicação serviço-a-serviço (leste-oeste); gateway cuida da entrada
(norte-sul).
\item \textbf{IaC (Infraestrutura como Código):} Terraform, Ansible ---
provisiona ambiente de forma declarativa e versionada.
\item \textbf{Kubernetes:} taints/tolerations, distribuições
RKE1/RKE2/K3s/EKS, Rancher.
\item \textbf{Armazenamento de objetos:} flat namespace, API REST.
\item \textbf{Fundamentos de SO que a FGV cobra como ``arquitetura'':}
paginação/memória virtual, E/S bloqueante, troca de contexto.
\end{itemize}

\begin{comosair}
Na dúvida entre duas alternativas, os gatilhos de distrator deste bloco
denunciam a errada: ``sempre'', ``exclusivamente'', ``elimina a necessidade
de'', ``idêntico ao''. Os pares que a FGV inverte aqui ---
vertical$\times$horizontal, REST$\times$SOAP, container$\times$VM,
fila$\times$tópico, web$\times$aplicação --- estão nas caixas vermelhas acima:
revise por elas, não por um resumo que as repita.
\end{comosair}
